\subsection{Validation Status}

Before reporting model performance, we validate that the benchmark
surface itself is reproducible. Table~\ref{tab:harness-validation}
summarises the offline run produced by
\path{benchmarks/scripts/run_benchmark_integration.py}. This run
checks benchmark unit tests, fixture-level integration tests, and
deterministic FullText/NaiveRAG smoke baselines without calling an LLM.
We therefore treat it as harness evidence, not as a headline
performance result.

\begin{table}[t]
    \centering
    \small
    \begin{tabular}{p{0.46\columnwidth}rll}
        \toprule
        Validation step & Cases/tests & Score/pass & Tokens \\
        \midrule
        Benchmark unit tests & 54 & 54 passed & -- \\
        Open-source fixture integration & 4 & 4 passed & -- \\
        LoCoMo FullText fixture & 3 & 1.0000 & 118 \\
        LoCoMo NaiveRAG fixture & 3 & 1.0000 & 118 \\
        LongMemEval FullText fixture & 4 & 0.7500 & 150 \\
        LongMemEval NaiveRAG fixture & 4 & 0.7500 & 150 \\
        \bottomrule
    \end{tabular}
    \caption{Offline harness validation run on 2026-06-24. These
    checks validate adapters, deterministic scorers, fixture wiring,
    and paper-table synchronization; they are not reported as live
    model performance.}
    \label{tab:harness-validation}
\end{table}


\subsection{Supplementary B0 Runtime Baseline}

The end-to-end B0 run completed all seven verifier checks in three of
five attempts, for an observed success rate of 60\%. Both rejected
attempts ended with \texttt{finish\_reason=length}; they remain in the
released ledger and in the effectiveness denominator. Across the three
accepted trajectories, mean input was 45,940 tokens, mean uncached input
was 13,428 tokens, the aggregate cached-token ratio was 70.8\%, and mean
wall time was 109.0 seconds. Table~\ref{tab:b0-wal-baseline} reports the
individual success-conditioned costs.

\begin{table*}[t]
    \centering
    \small
    \begin{tabular}{lrrrrrrr}
        \toprule
        Accepted run & Input & Cached & Uncached & Output & Model calls & Tool calls & Time (s) \\
        \midrule
        S1 & 40,080 & 30,208 & 9,872 & 2,211 & 13 & 12 & 97.5 \\
        S2 & 47,580 & 32,512 & 15,068 & 3,093 & 15 & 14 & 110.4 \\
        S3 & 50,161 & 34,816 & 15,345 & 4,755 & 10 & 9 & 119.3 \\
        \midrule
        Mean & 45,940 & 32,512 & 13,428 & 3,353 & 12.7 & 11.7 & 109.0 \\
        \bottomrule
    \end{tabular}
    \caption{Success-conditioned B0 full-replay costs on the
    Terminal-Bench 2.0 \texttt{db-wal-recovery} task. Five sequential
    LongCat-2.0 attempts used one frozen binary; three passed all seven
    verifier checks. The two rejected attempts are included in the
    separately reported 3/5 effectiveness rate, but excluded from this
    conditional efficiency table.}
    \label{tab:b0-wal-baseline}
\end{table*}


This small sample establishes a reproducible full-replay operating
point and exposes material trajectory variance, but it is neither a
confidence interval nor evidence that file-backed garbage collection
improves utility or cost. A treatment comparison requires the same
frozen binary, task, attempt count, and acceptance rule for FBGC, with
both all-attempt effectiveness and success-conditioned efficiency
reported.

\subsection{Main Live Results}

Table~\ref{tab:main} is reserved for live model runs produced by
\path{benchmarks/scripts/run_full_memory_benchmark.py}. The narrowed
claim requires the same sampled cases, same reader model, same
temperature, and same deterministic scorer for each method. Until the
live JSON artifacts are regenerated with \texttt{OPENAI\_\_*}
configured, the table intentionally remains empty rather than mixing
fixture checks with model-performance claims.

\begin{table}[t]
    \centering
    \small
    \begin{tabular}{p{0.48\columnwidth}ccc}
        \toprule
        Method & Accuracy & Tokens & Latency (s) \\
        \midrule
        FullText / flat window (B1) & -- & -- & -- \\
        NaiveRAG / top-$k$ retrieval (B2) & -- & -- & -- \\
        \midrule
        \textbf{\textsc{StructurePathMemory} (B3)} & \textbf{--} & \textbf{--} & \textbf{--} \\
        \bottomrule
    \end{tabular}
    \caption{Main live results on the narrowed benchmark set. Accuracy
    is the benchmark-native score; Tokens is mean prompt+completion
    tokens per trajectory; Latency is wall-clock seconds per
    trajectory.}
    \label{tab:main}
\end{table}

The performance discussion should be written only after the live
artifacts are present. The intended decision rule is conservative:
\textsc{Structure} supports the paper's claim only if it improves
tokens per scored point at comparable accuracy on at least two of the
three narrowed benchmarks.

\subsection{External Memory-Layer Baselines}

To make the comparison compatible with the LightMem/MemBase benchmark
suite, Table~\ref{tab:lightmem-locomo-baselines} records the published
LoCoMo baseline rows we use as external calibration data. The table is
generated from the repository script
\path{benchmarks/scripts/lightmem_baseline_report.py}, and CI
checks that the paper artifact stays in sync with the code-level
baseline registry.

% Generated by: python -m benchmarks.scripts.lightmem_baseline_report --format latex
\begin{table}[t]
    \centering
    \scriptsize
    \resizebox{\linewidth}{!}{%
    \begin{tabular}{lrrrrrr}
        \toprule
        Method & GPT judge & Qwen judge & Mem. Tok. & QA Tok. & Total Tok. & Runtime \\
        \midrule
        FullText & 73.83 & 73.18 & -- & 54{,}884.5 & 54{,}884.5 & 6{,}971 \\
        NaiveRAG & 63.64 & 63.12 & -- & 3{,}870.2 & 3{,}870.2 & 1{,}884 \\
        A-MEM & 64.16 & 60.71 & 11{,}494.3 & 10{,}170.6 & 21{,}664.9 & 67{,}084 \\
        MemoryOS(eval) & 58.25 & 61.04 & 2{,}870.0 & 7{,}649.3 & 10{,}519.4 & 26{,}129 \\
        MemoryOS(pypi) & 54.87 & 55.91 & 5{,}264.8 & 6{,}126.1 & 11{,}390.0 & 37{,}912 \\
        Mem0 & 36.49 & 37.01 & 24{,}304.9 & 1{,}488.6 & 25{,}793.5 & 120{,}175 \\
        Mem0(api) & 61.69 & 61.69 & 68{,}347.7 & 4{,}169.9 & 72{,}517.6 & 10{,}445 \\
        Mem0-g(api) & 60.32 & 59.48 & 69{,}684.8 & 4{,}389.1 & 74{,}074.0 & 10{,}926 \\
        \bottomrule
    \end{tabular}
    }
    \caption{External LoCoMo baseline data reported by LightMem/MemBase for a GPT-4o-mini backbone. Accuracy columns are percentages under GPT-4o-mini and Qwen2.5-32B judges. Token columns are in thousands. These rows are reference baselines for our evaluation plan, not new \textsc{Structure} measurements. Source: \citet{fang2026lightmem} and \url{https://github.com/zjunlp/LightMem\#experimental-results}.}
    \label{tab:lightmem-locomo-baselines}
\end{table}


\subsection{Hot-Research Benchmark Coverage}

The automated research scan identified LongMemEval-V2
\citep{wu2026longmemevalv2} and evidence-supported benchmark bounds
\citep{gao2026evidencebounds} as immediately actionable benchmark
gaps. We therefore added a fixture-level LongMemEval-V2 adapter that
preserves the benchmark's Insert/Query memory-system shape and scores
both answer correctness and supporting evidence identifiers. The new
coverage is reported in Table~\ref{tab:hot-research-coverage}; these
numbers are oracle fixture checks used to validate harness wiring, not
claims about full-dataset performance.

\begin{table}[t]
    \centering
    \small
    \begin{tabular}{lccc}
        \toprule
        Adapter & Local cases & Oracle score & Evidence-aware \\
        \midrule
        LongMemEval & 4 & 1.000 & no \\
        LoCoMo & 3 & 1.000 & no \\
        LongMemEval-V2 & 3 & 1.000 & yes \\
        \bottomrule
    \end{tabular}
    \caption{Hot-research benchmark coverage added by the
    LongMemEval-V2 adapter spike. Scores are fixture-level oracle
    checks from the shared benchmark runner, not full-dataset claims.}
    \label{tab:hot-research-coverage}
\end{table}


\subsection{Ablation Study}

The narrowed ablation focuses on mechanisms implemented in the
released benchmark path.
\textbf{(A1) Single level.} Every selected entry is materialised as
full \textsc{detail}; \textsc{glance} and \textsc{overview} no longer
reduce prompt cost.
\textbf{(A2) No rating.} The \texttt{rate\_context} feedback signal is
not used for filtering or ordering.
\textbf{(A3) No GC.} All run events remain visible to the replay path.
We do not ablate the audit log itself in the main paper because it is
part of the evidence contract used to reproduce runs.

\begin{table}[t]
    \centering
    \small
    \begin{tabular}{lcc}
        \toprule
        Configuration & Accuracy & Tokens \\
        \midrule
        Full system & -- & -- \\
        \midrule
        (A1) Single level & -- & -- \\
        (A2) No rating & -- & -- \\
        (A3) No GC & -- & -- \\
        \bottomrule
    \end{tabular}
    \caption{Narrowed ablation of implemented feedback and disclosure
    mechanisms.}
    \label{tab:ablation}
\end{table}

\subsection{Rating and GC Sensitivity}

Beyond on/off ablations, we sweep the two hyperparameters of the new
mechanisms. \textbf{Rating blend $\alpha$.} We vary
$\alpha \in \{0, 0.25, 0.5, 0.7, 0.9, 1.0\}$ and report the sweep only
if each point uses the same sampled cases and reader model.
\textbf{GC step-TTL.} We scale all non-pinned step-TTL values in the
default policy by $\beta \in \{0.5, 1.0, 2.0, 4.0, \infty\}$ and
interpret the curve as a cost--accuracy sensitivity analysis, not as a
claim that the hand-tuned default is globally optimal.

\subsection{Analysis}

\paragraph{Glance-to-detail ratio.} Live reports include selected and
available chunk counts; trajectory plots are generated from the
per-case JSONL artifacts when available.

\paragraph{Path locality.} We report mean prefix locality only for
\textsc{StructurePathMemory} runs whose per-case diagnostics include
accessed path prefixes.

\paragraph{Scalability with workspace size.} This scalability result
is outside the narrowed main claim and is reported as an appendix
diagnostic when the synthetic workspace-size sweep is run.

\paragraph{Qualitative example.} We include one annotated successful
run only after the corresponding live JSONL artifact is available, so
the narrative can be traced back to concrete retrieved chunks and
model outputs.
